\chapter{Implementation}
\label{ch:implementation}

The complete source code is publicly available at~\cite{chapidze2026opusdb}.

\section{Alignment with TL2}
This implementation follows the structure of Transactional Locking II
(TL2)~\cite{dice2006transactional}, but departs from the original in three
ways.

TL2 supports two lock-mapping policies: per-object (PO), where each reference
carries its own versioned lock, and per-stripe (PS), where each
reference hashes to a lock in a fixed array~\cite{dice2006transactional}.
Neither is used here. Ownership acquisition at write time instead goes
through a dedicated \texttt{java.lang.Object} monitor per reference, while
commit serialization is handled separately, by a single global commit lock.

TL2 discovers write conflicts at commit time; here they're caught earlier,
\emph{eagerly}, at the point of write, via \texttt{try-acquire!}, so a
losing transaction aborts on the spot rather than wasting work on a
transaction it's about to abort anyway.

TL2 also keeps only the most recently committed value of each reference.
This implementation adds \emph{Multi-Version Concurrency Control
  (MVCC)}~\cite{bernstein1983} on top, keeping a bounded version history
per reference so readers can find a consistent snapshot without
blocking writers (see Section~\ref{sec:mvcc-design}).

\section{Architecture Overview}
The \texttt{opusdb.atomic.stm} namespace defines four global data structures
and four constants that coordinate all transactional activity.

\begin{description}
  \item[\texttt{write-point}] An \texttt{AtomicLong} serving as the version
        clock, incremented atomically on every successful write-transaction
        commit to establish a total order over committed states.
  \item[\texttt{tx-counter}] A monotonically increasing \texttt{AtomicLong}
        that assigns unique transaction identifiers, whose ordering forms the
        basis of the ownership protocol.
  \item[\texttt{active-tx}] A \texttt{ConcurrentHashMap} mapping live
        transaction IDs to their transaction maps, used to locate and abort
        the owner of a conflicting reference.
  \item[\texttt{commit-lock}] A single global \texttt{Object} monitor that
        serializes read-set validation and write application across all write
        transactions.
\end{description}

Four constants tune the runtime behaviour:

\begin{description}
  \item[\texttt{max-history} (32)] Maximum size of a version history per
        reference; older entries are discarded on commit.
  \item[\texttt{max-retries} (96)] Maximum retries before throwing
        unconditionally. At roughly 2 seconds of cumulative delay, a
        still-failing transaction is assumed to be in livelock.
  \item[\texttt{base-sleep-ms} (1) and \texttt{max-sleep-ms} (32)] Bounds
        of the exponential backoff delay applied after the fourth retry,
        capped at 32\,ms to keep worst-case latency bounded.
\end{description}

The current transaction is bound to the thread-local variable \texttt{*tx*},
with \texttt{nil} indicating no active transaction context.

\section{Reference and Version History}

A transactional reference is created by \texttt{ref}, which packages the
initial value, version history, and ownership metadata into a single Clojure
\texttt{atom}:

\begin{listing}[H]
  \begin{minted}{clojure}
(defn ref [initial-val]
  (let [wp    (.get write-point)
        entry {:value initial-val :write-point wp}]
    (atom {:value       initial-val
           :write-point wp
           :history     [entry]
           :owner       nil
           :lock        (Object.)})))
\end{minted}
  \caption{Reference structure.}
  \label{lst:ref-structure}
\end{listing}

The \texttt{:history} vector stores committed versions as
\texttt{\{:value v :write-point w\}} entries. The \texttt{:lock} field is a
dedicated \texttt{Object} monitor; \texttt{try-acquire!} synchronizes on it
when competing transactions claim ownership of the same reference.

Version history is bounded by \texttt{max-history = 32}. When a commit would
push the count beyond this limit, \texttt{trim-history} drops the oldest
entry before appending the new one:

\begin{listing}[H]
  \begin{minted}{clojure}
(defn- trim-history [history]
  (if (> (count history) max-history)
    (vec (take-last max-history history))
    history))
\end{minted}
  \caption{History trimming.}
  \label{lst:trim-history}
\end{listing}

\section{Transaction Lifecycle}

\begin{figure}[H]
  \centering
  \begin{tikzpicture}[
    font=\sffamily\small,
    every node/.style={align=center},
    N/.style={
        rectangle, rounded corners=3pt,
        draw=black!45, fill=gray!8,
        minimum width=3.4cm, minimum height=0.72cm,
        text width=3.2cm},
    D/.style={N, draw=blue!50!black, fill=blue!7},
    OK/.style={N, draw=teal!55!black, fill=teal!8},
    AB/.style={N, draw=red!48!black,  fill=red!7},
    T/.style={N, draw=black!55, fill=gray!16,
        font=\sffamily\small\bfseries},
    a/.style={-{Stealth[scale=0.75]}, thick},
    d/.style={-{Stealth[scale=0.75]}, thick, dashed, draw=gray!55},
    ]
    \def\lx{-4.2}
    \def\cx{0}
    \def\rx{4.2}
    \def\loopx{-7.0}
    \def\routeR{7.4}
    \def\readonlyL{-2.2}

    \node[T]  (s)  at (\cx,   0.0) {dosync};
    \node[N]  (mt) at (\cx,  -1.8) {make-tx\\
    {\scriptsize snapshot write-point,\\register in active-tx}};
    \node[N]  (bf) at (\cx,  -4.0) {execute body-fn\\
    {\scriptsize check-active! per op\\
    try-acquire! on write\\
    build read/write sets}};
    \node[D]  (oc) at (\cx,  -6.2) {outcome?};
    \node[D]  (we) at (\cx,  -8.0) {write-set empty?};

    \node[N]  (al) at (\rx,  -9.8) {acquire commit-lock\\
    {\scriptsize re-check status}};
    \node[N]  (vr) at (\rx, -11.8) {validate-reads!\\
    {\scriptsize versions $\leq$ write-point?}};
    \node[N]  (aw) at (\rx, -13.8) {apply-writes!\\
    {\scriptsize increment write-point,\\stamp + append versions}};
    \node[N]  (rl) at (\rx, -15.6) {release commit-lock};
    \node[OK] (cm) at (\cx, -17.4) {tx/commit!\\
    {\scriptsize dispatch on-commit hooks}};

    \node[AB] (rb) at (\lx,  -6.2) {tx/rollback!\\
    {\scriptsize dispatch on-rollback hooks}};
    \node[N]  (bk) at (\lx,  -8.6) {backoff\\
    {\scriptsize exponential + jitter}};
    \node[D]  (mr) at (\lx, -10.4) {retries $\geq$ 96?};
    \node[AB] (th) at (\lx, -12.2) {throw\\
    {\scriptsize "STM exceeded max retries"}};

    \draw[a] (s)  -- (mt);
    \draw[a] (mt) -- (bf);
    \draw[a] (bf) -- (oc);
    \draw[a] (oc) -- node[right, font=\scriptsize\itshape]{ok} (we);

    \draw[a] (we.west)
    -- node[above, font=\scriptsize\itshape]{yes}
    (\readonlyL,-8.0)
    -- (\readonlyL,-17.4)
    -- (cm.west);

    \draw[a] (we.east)
    -- node[above, font=\scriptsize\itshape]{no}
    (\rx,-8.0)
    -- (al.north);
    \draw[a] (al) -- (vr);
    \draw[a] (vr) -- node[right, font=\scriptsize\itshape]{pass} (aw);
    \draw[a] (aw) -- (rl);
    \draw[a] (rl.south) -- (\rx,-16.5) -- (cm.east);

    \draw[d] (vr.east)
    -- node[right, font=\scriptsize\itshape]{fail}
    (\routeR,-11.8)
    -- (\routeR,-6.2)
    -- (oc.east);

    \draw[a] (oc.west)
    -- node[above, font=\scriptsize\itshape]{abort}
    (rb.east);
    \draw[a] (rb) -- (bk);
    \draw[a] (bk) -- (mr);
    \draw[a] (mr) -- node[right, font=\scriptsize\itshape]{yes} (th);

    \draw[a] (mr.west)
    -- node[above, font=\scriptsize\itshape]{no}
    (\loopx,-10.4)
    -- (\loopx,-4.0)
    -- (bf.west);

  \end{tikzpicture}
  \caption{Transaction execution flow. The speculative phase builds read/write
    sets, checking status and acquiring ownership on each write; on success
    the commit path validates read-set, applies writes under the commit lock
    (read-only transactions bypass the commit lock entirely), and dispatches
    commit hooks; on failure the abort path dispatches rollback hooks, backs off,
    and retries with a fresh snapshot.}
  \label{fig:tx-flow}
\end{figure}

The subsections below walk through each phase.

\subsection{Creation}

Each transaction is created by \texttt{make-tx}, which atomically increments
\texttt{tx-counter} to obtain a unique identifier, snapshots the current
\texttt{write-point} as the read point, initializes empty read and write sets,
and registers the transaction in \texttt{active-tx}:

\begin{listing}[H]
  \begin{minted}{clojure}
(defn- make-tx [retry-count]
  (let [id (.incrementAndGet tx-counter)
        tx {:id          id
            :read-point  (.get write-point)
            :read-set    (IdentityHashMap.)
            :write-set   (IdentityHashMap.)
            :retry-count retry-count
            :status      (volatile! ::active)}]
    (.put active-tx id tx)
    tx))
\end{minted}
  \caption{Transaction creation. The read point is snapshotted at creation
    time, defining the consistent view of memory this transaction will observe.}
  \label{lst:make-tx}
\end{listing}

Both the read-set and write-set use \texttt{IdentityHashMap}, which compares
keys by reference identity rather than value equality. This matters because
transactional references are \texttt{atom} objects: their identity stays
fixed for life, even as their contents change on every commit.

The read-set also caches the \texttt{:write-point} observed at read time
alongside the value. Since the reference's own \texttt{:write-point} field
gets overwritten by later commits, read-set validation would otherwise have
nothing left to compare against when checking for conflicts.

\subsection{Reading}
\label{sec:reading}

\texttt{deref} resolves a reference value through three ordered lookups:

\begin{listing}[H]
  \begin{minted}{clojure}
(defn deref [target-ref]
  (if-not *tx*
    (:value @target-ref)
    (let [tx        *tx*
          write-set (:write-set tx)
          read-set  (:read-set tx)]
      (check-active! tx)

      (cond
        (.containsKey write-set target-ref)
        (.get write-set target-ref)

        (.containsKey read-set target-ref)
        (:value (.get read-set target-ref))

        :else
        (let [snapshot (find-entry (:read-point tx)
                                   (:history @target-ref))]
          (when-not snapshot (abort))
          (.put read-set target-ref snapshot)
          (:value snapshot))))))
  \end{minted}
  \caption{Transactional read. The write-set is checked first for
    read-your-writes semantics; the read-set returns the previously cached
    snapshot; the version history is checked last.}
  \label{lst:deref}
\end{listing}

Checking the write-set first guarantees read-your-writes: a transaction
always sees whatever it has already buffered. After that comes the
read-set: if the reference was already read earlier in the same
transaction, its cached snapshot comes back directly. Only when neither set
has it does \texttt{find-entry} search the history, returning the latest
committed entry whose write point doesn't exceed the transaction's read
point. If the history has been trimmed past the point where such an entry
exists, the transaction aborts and retries with a fresh snapshot.

\subsection{Writing}

\texttt{ref-set} buffers a write into the transaction's write-set, but only
after acquiring ownership of the target reference via \texttt{try-acquire!}:

\begin{listing}[H]
  \begin{minted}{clojure}
(defn ref-set [target-ref new-val]
  (when-not *tx*
    (throw (IllegalStateException. "ref-set outside transaction")))
  (let [tx *tx*]
    (check-active! tx)
    (when-not (try-acquire! target-ref tx)
      (abort))
    (.put ^IdentityHashMap (:write-set tx) target-ref new-val)
    new-val))
\end{minted}
  \caption{Transactional write. Ownership must be acquired before buffering;
    failure to acquire triggers an immediate abort.}
  \label{lst:ref-set}
\end{listing}

\texttt{alter} is a thin convenience wrapper: it reads the current value and
hands it to a caller-supplied function before delegating to \texttt{ref-set}:

\begin{listing}[H]
  \begin{minted}{clojure}
(defn alter [target-ref f & args]
  (ref-set target-ref (apply f (deref target-ref) args)))
\end{minted}
  \caption{alter reads the current value and writes the result of applying
    \texttt{f}, reusing the ownership and buffering logic of \texttt{ref-set}.}
  \label{lst:alter}
\end{listing}

\subsection{Committing}

\texttt{commit!} finalizes a transaction. Read-only transactions, those
with an empty write-set, skip the global commit lock entirely and commit
right away, since the version history already guarantees their snapshot is
consistent. Write transactions don't get that shortcut: they enter the
commit lock, re-check their status (ownership theft can happen while
waiting), validate the read-set, and only then apply their writes:

\begin{listing}[H]
  \begin{minted}{clojure}
(defn- commit! [tx]
  (check-active! tx)
  (let [{:keys [id read-set write-set status]} tx]
    (when (seq write-set)
      (locking commit-lock
        (check-active! tx)
        (validate-reads! read-set)
        (apply-writes! write-set)))
    (vreset! status ::committed)
    (.remove active-tx id)
    (tx/commit! id)))
\end{minted}
  \caption{Commit protocol. Read-only transactions skip the commit lock;
    write transactions validate reads before applying writes.}
  \label{lst:commit}
\end{listing}

\texttt{apply-writes!} increments \texttt{write-point} exactly once and uses
that value to timestamp every write in the batch, which is what keeps all
writes from one commit atomic with respect to readers:

\begin{listing}[H]
  \begin{minted}{clojure}
(defn- apply-writes! [^IdentityHashMap write-set]
  (let [commit-point (.incrementAndGet write-point)]
    (doseq [[r new-val] write-set]
      (swap! r (fn [{:keys [history] :as state}]
                 (assoc state
                        :value       new-val
                        :write-point commit-point
                        :owner       nil
                        :history     (trim-history
                                      (conj history
                                            {:value       new-val
                                             :write-point commit-point}))))))))
\end{minted}
  \caption{Write application. A single increment of the global write point
    timestamps all writes in a commit atomically.}
  \label{lst:apply-writes}
\end{listing}

\subsection{Retry and Backoff}
\label{sec:retry}

When a transaction aborts, \texttt{next-tx} dispatches its rollback hooks,
removes it from \texttt{active-tx}, and hands back a fresh transaction with
an incremented retry count. The first four retries carry no sleep penalty at
at all, since short-lived conflicts that resolve quickly shouldn't have to pay for a
\texttt{Thread/sleep} call. Past the fourth retry, backoff kicks in, and if
\texttt{max-retries} $= 96$ is ever reached the system throws unconditionally,
on the assumption that this is most likely a livelock by that point.

For retry count $n > 4$, the sleep duration is:
\begin{equation}
  S = b(n)\!\left(1 + \tfrac{1}{4}X\right), \quad X \sim \mathcal{U}[-1,1]
\end{equation}
with bounds $S_{\min} = \tfrac{3}{4}b(n)$, $S_{\max} = \tfrac{5}{4}b(n)$,
and base delay:
\begin{equation}
  b(n) = \min\!\left(32,\; 2^{\min(n-4,\,6)}\right)\,\text{ms}
\end{equation}
giving $b(5) = 2\,\text{ms}$, $b(8) = 16\,\text{ms}$, and $b(n) = 32\,\text{ms}$
for all $n \geq 9$. The $\pm 25\%$ jitter exists to prevent the
thundering-herd problem\footnote{The thundering-herd problem occurs when multiple transactions
  abort simultaneously and retry at the same time, causing them to conflict again
  and repeat the cycle. Jitter breaks the symmetry by randomizing each
  transaction's retry delay.}, spreading retry attempts out over time; the
$32\,\text{ms}$ cap keeps worst-case per-retry latency bounded no matter how
severe the contention gets.

\begin{listing}[H]
  \begin{minted}{clojure}
(defn- next-tx [current]
  (tx/rollback! (:id current))
  (.remove active-tx (:id current))
  (let [retries (:retry-count current)]
    (when (>= retries max-retries)
      (throw (ex-info "STM exceeded max retries" {:retries retries})))
    (when (> retries 4)
      (let [base   (min max-sleep-ms
                        (bit-shift-left base-sleep-ms
                                        (min (- retries 4) 6)))
            jitter (long (* base 0.25
                            (- (* 2 (.nextDouble
                                      (ThreadLocalRandom/current)))
                               1)))]
        (Thread/sleep (max 0 (long (+ base jitter))))))
    (make-tx (inc retries))))
\end{minted}
  \caption{Retry with bounded exponential backoff. The four-retry period
    avoids sleep overhead for short-lived conflicts.}
  \label{lst:next-tx}
\end{listing}

\subsection{The dosync Macro}

\texttt{dosync} is a thin macro over \texttt{sync}, which checks whether a
transaction is already active on the current thread; if not, \texttt{run-tx}
starts a new transaction loop:

\begin{listing}[H]
  \begin{minted}{clojure}
(defn sync [body-fn]
  (if *tx*
    (body-fn)
    (run-tx body-fn)))

(defmacro dosync [& body]
  `(sync (fn* [] ~@body)))
\end{minted}
  \caption{\texttt{dosync} delegates to \texttt{sync}, which joins an existing
    transaction if one is active, or starts a new one via \texttt{run-tx}.}
  \label{lst:dosync}
\end{listing}

\paragraph{Nested transaction semantics.}
When \texttt{dosync} runs on a thread that already has a transaction bound
to \texttt{*tx*}, \texttt{sync} just invokes the body inside that existing
transaction rather than opening a new one, which is \emph{flat nesting}. The inner
call shares the outer transaction's read-set, write-set, and snapshot, so an
abort anywhere in the inner body propagates outward and the whole
transaction retries. Hooks registered inside the inner body get associated
with the outer transaction's identifier, so they still fire at the right
lifecycle boundary. This is exactly the composability property described in
Section~\ref{sec:intro}: \texttt{withdraw} and \texttt{deposit} each have
their own \texttt{dosync}, but called inside an enclosing \texttt{dosync}
they end up participating in a single atomic transaction.
Figure~\ref{fig:tx-flow} lays out the full execution flow.

\section{Conflict Resolution}
\label{sec:conflict-resolution}

Rather than waiting until commit time to detect conflicts, the system
resolves write conflicts the moment a transaction first touches a reference.
\texttt{try-acquire!} claims ownership under that reference's per-ref
monitor, following an ownership protocol that favours the younger
transaction: a transaction with a higher identifier preempts whoever holds
ownership currently, marking them aborted, which they'll notice on
their next \texttt{check-active!} call. An older transaction that loses the
contest aborts itself right away and retries. The \texttt{max-retries} bound
from Section~\ref{sec:retry} keeps a repeatedly preempted transaction from
spinning forever.

\begin{listing}[H]
  \begin{minted}{clojure}
(defn- try-acquire! [target-ref tx]
  (let [claimer-id (:id tx)
        lock       (:lock @target-ref)]
    (monitor-enter lock)
    (try
      (let [owner-id (:owner @target-ref)]
        (cond
          (nil? owner-id)
          (do (swap! target-ref assoc :owner claimer-id) true)
          (= owner-id claimer-id)
          true
          (> claimer-id owner-id)
          (do (when-some [owner-tx (.get active-tx owner-id)]
                (when (= @(:status owner-tx) ::active)
                  (vreset! (:status owner-tx) ::aborted)))
              (swap! target-ref assoc :owner claimer-id)
              true)
          :else false))
      (finally
        (monitor-exit lock)))))
\end{minted}
  \caption{Ownership acquisition. Younger transactions preempt the older
    owner, which aborts and retries.}
  \label{lst:try-acquire}
\end{listing}

The four \texttt{cond} branches cover every possible ownership state. An
unowned reference gets claimed immediately. A reference already owned by the
same transaction is just a no-op, which comes up when \texttt{ref-set} gets
called more than once on the same ref. If the current owner is older, the
claimer marks it \textsc{Aborted} in \texttt{active-tx} atomically and takes
over; if the current owner is younger, the claimer backs off instead,
returning \texttt{false} and aborting so the younger transaction has room to
finish.

The aborted owner finds out about its own invalidation the next time
\texttt{check-active!} runs, which happens before every critical operation,
and exits without committing any partial state.

\section{MVCC Design}
\label{sec:mvcc-design}

The system extends TL2 with multi-version concurrency
control~\cite{bernstein1983} by maintaining a bounded version history per
reference. Each history entry records a committed value and the write point
it was committed at. When a transaction commits at write point $w_c$, the
history of each written reference updates as:
\[
  \text{hist}_{\text{new}}(r) = \text{trim}_{32}\!\left(
  \text{hist}_{\text{old}}(r) \cup \bigl\{ (w_c,\; WS[r]) \bigr\}
  \right)
\]
where $\text{trim}_{32}$ drops the oldest entry if the count would exceed
\texttt{max-history}.

A transaction with read point $w_r$ retrieves the visible value of reference
$r$ via \texttt{find-entry}, which walks the history newest-to-oldest and
returns the first entry whose write point doesn't exceed $w_r$:

\begin{listing}[H]
  \begin{minted}{clojure}
(defn- find-entry [read-point history]
  (loop [i (dec (count history))]
    (when (>= i 0)
      (let [entry (nth history i)]
        (if (<= (:write-point entry) read-point)
          entry
          (recur (dec i)))))))
\end{minted}
  \caption{Version history scan. Newest-first traversal returns the latest
    version visible at the transaction's read point.}
  \label{lst:find-entry}
\end{listing}

In practice, this means readers and writers never block each other. A reader
at $w_r$ simply picks up a snapshot from before any concurrent writer's
commit, so neither side ever waits on the other. It's this property that
drives the strong scaling seen in the read-heavy benchmark of
Section~\ref{sec:read-write-heavy}: readers never touch the commit-lock
path at all.

If a transaction's read point falls below the oldest retained entry, no
valid snapshot exists, and the transaction aborts immediately and retries
with a fresh read point.

\section{Transactional Event System}

Side-effects are unsafe inside a transaction body, since that body may run
many times before a commit actually succeeds. Logging, sending a
notification, or writing to an external system on the first attempt would
fire again on every subsequent retry. The transactional event system in the
\texttt{opusdb.atomic.tx-hooks} namespace solves this by deferring handler
execution to a well-defined lifecycle boundary instead.

Two hooks are provided. \texttt{on-commit} schedules a handler to run
exactly once, after a successful commit. \texttt{on-rollback} schedules one
to run on each abort, before the next retry, which is useful for cleanup or
diagnostic logging:

\begin{listing}[H]
  \begin{minted}{clojure}
(defn on-commit [handler]
  (if *tx*
    (tx/on-commit (:id *tx*) handler)
    (throw (IllegalStateException. "on-commit outside transaction"))))

(defn on-rollback [handler]
  (if *tx*
    (tx/on-rollback (:id *tx*) handler)
    (throw (IllegalStateException. "on-rollback outside transaction"))))
\end{minted}
  \caption{Hook registration. Handlers are associated with a transaction ID
    and dispatched at the appropriate lifecycle event.}
  \label{lst:hook-registration}
\end{listing}

Storage underneath uses two \texttt{ConcurrentHashMap} instances, one for
commit handlers and one for rollback handlers, both indexed by transaction
ID:

\begin{listing}[H]
  \begin{minted}{clojure}
(def ^:private ^ConcurrentHashMap COMMIT_EVENTS   (ConcurrentHashMap.))
(def ^:private ^ConcurrentHashMap ROLLBACK_EVENTS (ConcurrentHashMap.))

(defn on-commit [id fun]
  (.compute COMMIT_EVENTS id
            (fn [_ fns] (conj (or fns []) fun))))

(defn on-rollback [id fun]
  (.compute ROLLBACK_EVENTS id
            (fn [_ fns] (conj (or fns []) fun))))
\end{minted}
  \caption{Core event tables and handler accumulation.}
  \label{lst:event-tables}
\end{listing}

\texttt{ConcurrentHashMap.compute} is atomic, so accumulating handlers needs
no extra locking on top. Dispatch happens at the commit and rollback
boundaries; on a successful commit, the rollback handlers for that
transaction get discarded too, since they'll never be needed:

\begin{listing}[H]
  \begin{minted}{clojure}
(defn commit! [id]
  (when-let [fns (.remove COMMIT_EVENTS id)]
    (doseq [f fns] (f))
    (.remove ROLLBACK_EVENTS id)))

(defn rollback! [id]
  (when-let [fns (.remove ROLLBACK_EVENTS id)]
    (doseq [f fns] (f))))
\end{minted}
  \caption{Dispatch at commit and rollback boundaries.}
  \label{lst:dispatch}
\end{listing}